\documentclass[10pt,a4paper]{amsart}
\usepackage[margin=19mm]{geometry}
\usepackage{amsmath,amssymb,mathtools}
\usepackage[scr=rsfs,cal=euler]{mathalfa}
\usepackage{xfrac}
\usepackage[unicode,hidelinks,bookmarks=false]{hyperref}
\newcommand{\ie}{i.e.\ }
\newcommand{\cf}{cf.\ }
\newcommand{\resp}{respectively\ }
\newcommand{\Subsection}{\subsection}
\providecommand{\nobreakdash}{\nobreak\hbox{-}}
\setlength{\emergencystretch}{2em}
\multlinegap0pt
\usepackage{bm}
\usepackage{bbm}
\usepackage{dsfont}
\usepackage{xspace}
\usepackage{cancel}
\usepackage{mathtools}
\usepackage{amsthm}
\makeatletter
\@ifundefined{theorem}{\newtheorem{theorem}{Theorem}}{}
\@ifundefined{lemma}{\newtheorem{lemma}[theorem]{Lemma}}{}
\@ifundefined{proposition}{\newtheorem{proposition}[theorem]{Proposition}}{}
\makeatother
\newcommand \G {\mathbb{G}}
\newcommand \N {\mathbb{N}}
\newcommand \auto {{\rm Aut}}
\newcommand \lc {{\rm lc}}
\renewcommand \a {A}
\renewcommand \k {K}
\title[The isotropy group of a derivation on a Danielewski-type alg]{The isotropy group of a derivation on a Danielewski-type algebra}
\author{Abdessamad Ahouita, Rene Baltazar, M'hammed El Kahoui, Sergey Gaifullin}
\date{Source: arXiv:2510.07059v2 (CC BY 4.0)}
\begin{document}
\maketitle

\noindent\emph{Source and licence.} The isotropy group of a derivation on a Danielewski-type algebra. Version arXiv:2510.07059v2 (CC BY 4.0), distributed under the Creative Commons Attribution 4.0 International licence, as recorded in the arXiv licence metadata for this exact version and shown on the abstract page https://arxiv.org/abs/2510.07059v2. Full rights notice: licenses/arxiv-2510.07059-CC-BY-4.0.txt.\par\smallskip

\noindent\textit{Editorial scope.} Counted unit: the Lemma whose source label is \texttt{closedSubgroup} in the article (source0.tex lines 234--235, source label \texttt{closedSubgroup}), delivered together with the article's own complete original proof of it (source0.tex lines 236--279, one \texttt{proof} environment of 44 lines). No mathematics is written, repaired or re-derived: statement and proof are the article's own text line for line. Required context retained uncounted: xbasisProposition (lines 146--148). Every definition, display and label of those lines is retained unchanged. Omitted from this package and not redistributed: 1--145, 149--233, 280--674. Nothing inside a retained excerpt is omitted or altered: every retained body is delivered exactly as the article's lines are written, so no omission receipt of any kind accompanies this package. Citations: the retained excerpts carry no citation command at all, so no bibliography entry of the article is transcribed and no reference is redistributed. Reference adaptation: none was needed and none was made; every reference inside the delivered unit resolves inside it, so no unresolved reference or citation is produced. Printed slips kept literal: no printed word, symbol, hypothesis or number of the counted statement or of its proof was altered. Layout and class: the article's own class is \texttt{article} with packages including graphicx, tikz, babel, geometry, amsmath, amssymb, latexsym, hyperref, amsfonts, amsthm, color, indentfirst, multirow, mathrsfs; the wrapper is the lane's pinned \texttt{amsart} editorial wrapper at 10pt on A4, which restates the theorem-like environments the retained text needs and adds the article's own macro definitions that the retained text needs (\texttt{\textbackslash G}, \texttt{\textbackslash N}, \texttt{\textbackslash a}, \texttt{\textbackslash auto}, \texttt{\textbackslash k}, \texttt{\textbackslash lc}), with extra wrapper packages (bm, bbm, dsfont, xspace, cancel, mathtools). The delivered numbering is the unit's own. Assets: no retained span contains a figure environment, a graphics-inclusion command, a PDF-page inclusion, a table or a TikZ picture; no image file is copied into this package, the package declares no asset dependency and no image file is redistributed with it. Licence: version arXiv:2510.07059v2 of this article is distributed under the Creative Commons Attribution 4.0 International licence, as recorded in the arXiv licence metadata for this exact version and shown on the abstract page https://arxiv.org/abs/2510.07059v2. A full rights notice is delivered at licenses/arxiv-2510.07059-CC-BY-4.0.txt. Characters: every retained excerpt is pure ASCII, so the delivered engine, XeLaTeX with its default Latin Modern text font, reports no missing character.
\begin{proposition}\label{xbasisProposition}Let $c\in\k[x]$ be a non-constant polynomial and $q(x,y)\in\k[x,y]$ be a quasi-monic polynomial of degree $d\geq 2$ with respect to $y$. Then 
	$$\big\{ \bar{y}^i\bar{z}^j\mid 0\leq i<d, \;j\in\N\}$$ is a $\xi_{c,q}$-basis of $\a_{c,q}$ over $\k[\bar{x}]=\ker(\xi_{c,q})$.
\end{proposition}

\begin{lemma}\label{closedSubgroup}Let $\delta$ be a $\k$-derivation of $\a_{c,q}$. Then $\auto_{\k}(\a_{c,q},\delta)$ is a closed ind-subgroup of $\auto_{\k}(\a_{c,q})$. If moreover $\delta$ is not locally nilpotent then $\auto_{\k}(\a_{c,q},\delta)$ is an algebraic subgroup of $\auto_{\k}(\a_{c,q})$.
\end{lemma}

\begin{proof}In order to prove that $\auto_{\k}(\a_{c,q},\delta)$ is a closed subgroup of $\auto_{\k}(\a_{c,q})$ we need to show that $\auto_{\k}(\a_{c,q},\delta)\cap \mathcal{G}_m$ is closed in $\mathcal{G}_m$ for every $m\geq 0$, where $$\mathcal{G}_m=\big\{ \sigma\in\auto_{\k}(\a_{c,q})\mid \deg(\ell)\leq m\big\},$$ and $\ell(x)\in\k[x]$ is such that $\sigma(\bar{y})=u\bar{y}+c(\bar{x})\ell(\bar{x})$.
	
	\medskip Let $k\geq 0$ be such that \begin{equation}\label{cancellationOfZ}c^{k}(\bar{x})\delta(\bar{x})=f_1(\bar{x},\bar{y}) \quad \text{and} \quad c^{k}(\bar{x})\delta(\bar{y})=f_2(\bar{x},\bar{y}),\end{equation} where $f_i\in \k[x,y]=K^{[2]}$ depends only on $\delta$. Let $\sigma\in\mathcal{G}_m$ and write $\sigma(\bar{x})=e\bar{x}$ and $\sigma(\bar{y})=u\bar{y}+h(\bar{x})$, where $(e,u)\in\G_{c,q}$ and $h(x)=c(x)\ell(x)$ and $\deg(\ell)\leq m$. Then by applying $\sigma$ to the equalities in (\ref{cancellationOfZ}), and taking into account $c(ex)=e^nc(x)$, we have $$\begin{array}{ccc}f_1(e\bar{x},u\bar{y}+h(\bar{x}))&=&
		e^{kn}c^k(\bar{x})\sigma\delta(\bar{x}),\\
		f_2(e\bar{x},u\bar{y}+h(\bar{x}))&=&e^{kn}c^k(\bar{x})\sigma\delta(\bar{y}).\end{array}$$
	On the other hand, we have
	$$\begin{array}{ccc}
		e^{kn}c^k(\bar{x})\delta\sigma(\bar{x})&=&e^{kn+1}f_1(\bar{x},\bar{y}),\\
		e^{kn}c^k(\bar{x})\delta\sigma(\bar{y})&=&e^{kn}f_2(\bar{x},\bar{y})+e^{kn}h^{\prime}(\bar{x})f_1(\bar{x},\bar{y}).
	\end{array}$$
	Since $\a_{c,q}$ is an integral domain and $c(\bar{x})\bar{z}=q(\bar{x},\bar{y})$ it follows that $\sigma\delta=\delta\sigma$ if and only if
	\begin{equation}\label{firstCondition}f_1(e\bar{x},u\bar{y}+h(\bar{x}))-e^{kn+1}f_1(\bar{x},\bar{y})=0
	\end{equation}
	and
	\begin{equation}\label{secondCondition}f_2(e\bar{x},u\bar{y}+h(\bar{x}))-e^{kn}uf_2(\bar{x},\bar{y})-e^{kn}h^{\prime}(\bar{x})f_1(\bar{x},\bar{y})=0.
	\end{equation}
	Since moreover $\k[\bar{x},\bar{y}]=\k^{[2]}$ the equations (\ref{firstCondition}) and (\ref{secondCondition}) yield a system of polynomial equations in terms of $e,u$ and the coefficients of $\ell(x)$, where $h(x)=c(x)\ell(x)$. This proves that $\auto_{\k}(\a_{c,q},\delta)\cap \mathcal{G}_m$ is closed in $\mathcal{G}_m$, and since this holds for every $m\geq 0$ it follows that $\auto_{\k}(\a_{c,q},\delta)$ is a closed ind-subgroup of $\auto_{\k}(\a_{c,q})$.
	
	\medskip Assume now that $\delta$ is not locally nilpotent. In order to prove that $\auto_{\k}(\a_{c,q},\delta)$ is an algebraic subgroup of $\auto_{\k}(\a_{c,q})$ we only need to show that $\auto_{\k}(\a_{c,q},\delta)\subseteq \mathcal{G}_m$ for some positive integer $m$. That is, for every $\sigma\in\auto_{\k}(\a_{c,q},\delta)$, with $\sigma(\bar{y})=u\bar{y}+h(\bar{x})$, we have $\deg(h)\leq m+n$ where $n=\deg(c)$.
	
	\medskip Let $\sigma\in\auto_{\k}(\a_{c,q},\delta)$, and let $(e,u)\in \G_{c,q}$ and $h(x)\in c(x)\k[x]$ be such that $$\sigma(\bar{x})=e\bar{x},\quad \sigma(\bar{y})=u\bar{y}+h(\bar{x}).$$ Then we have several possibilities, depending on the nature of $f_1$ and $f_2$.
	
	\medskip -- The case $f_1(x,y)\in K[x,y] \setminus \k[x]$. From the equation (\ref{firstCondition}) we have
	$$	f_1(e\bar{x},u\bar{y}+h(\bar{x}))= e^{kn+1} f_1(\bar{x}, \bar{y})$$
	and since $\k[\bar{x},\bar{y}]=K^{[2]}$ we actually have \begin{equation}\label{conditioncommefromxx} f_1(ex,uy+h(x))=e^{kn+1}f_1(x,y).\end{equation}
	Let $s = \deg_y(f_1) \geq 1$ and let us write $$ f_1(x, y) = \sum_{i=0}^{s} f_{1,i}(x) y^i,$$ where $ f_{1,i}(x) \in \k[x]$.
	By looking at the coefficient of  degree $s-1$ with respect to $y$ in both sides of \eqref{conditioncommefromxx}
	we obtain $$u^{s-1} f_{1,s-1}(ex)+su^{s-1} f_{1,s}(ex)h(x)=e^{kn+1}f_{1,s-1}(x).$$ Since $f_{1,s}(x)\neq0$ we have $\deg(h)\leq \deg_{x}(f_{1,s-1})\leq \deg_{x}(f_{1})$.
	
	\medskip -- The case $f_1(x,y)=f_1(x)\in\k[x]$ and $s=\deg_y(f_2)\geq 2$. From the equation (\ref{secondCondition}) we have 
	$$f_2(e\bar{x},u\bar{y}+h(\bar{x}))=e^{kn}u f_2(\bar{x},\bar{y}) + e^{kn}h^{\prime}(\bar{x})f_1(\bar{x})$$
	and since $\k[\bar{x},\bar{y}]=\k^{[2]}$ we actually have \begin{equation}\label{conditioncommefromyy} f_2(ex,uy+h(x))=e^{kn}u f_2(x, y)+e^{kn} h^{\prime}(x)f_1(x).\end{equation}
	Let us write $$ f_2(x, y) = \sum_{i=0}^{s} f_{2,i}(x) y^i,
	$$ where $ f_{2,i}(x) \in \k[x]$. By looking at the coefficient of  degree $s-1\geq 1$ with respect to $y$ in both sides of \eqref{conditioncommefromyy}
	we obtain $$u^{s-1} f_{2,s-1}(ex)+su^{s-1} f_{2,s}(ex)h(x)=e^{kn}uf_{2,s-1}(x).$$ Since $f_{2,s}(x)\neq0$, we have $\deg(h)\leq \deg_{x}(f_{2,s-1})\leq \deg_{x}(f_{2})$.
	
	\medskip -- The case $f_1(x,y)=f_1(x)\in\k[x]$ and $s=\deg_y(f_2)=1$. To simplify, we let $\xi$ stand for $\xi_{c,q}$ and recall that $\deg_{\xi}$ denotes the degree function associated to $\xi$. The assumption that $\deg_y(f_2)=1$ implies that $\deg_{\xi}(\delta(\bar{y}))=1$ and from Proposition \ref{xbasisProposition} it follows that \begin{equation}\label{equationThirdCase}\delta(\bar{y})=a(\bar{x})\bar{y}+b(\bar{x})\end{equation} where $a(x),b(x)\in\k[x]$ and $a(x)\neq 0$. On the other hand, since $c(\bar{x})^k\delta(\bar{x})=f_1(\bar{x})\in\k[\bar{x}]$ and $c(\bar{x})\neq 0$ and $\k[\bar{x}]$ is factorially closed in $\a_{c,q}$ it follows that $\delta(\bar{x})=g(\bar{x})\in\k[\bar{x}]$.
	
	\medskip By applying $\sigma$ to both sides of (\ref{equationThirdCase}), and taking into account $\sigma\delta=\delta\sigma$ and Proposition \ref{xbasisProposition} and $\k[\bar{x},\bar{y}]=\k^{[2]}$ we obtain $a(ex)=a(x)$ and \begin{equation}\label{equationThirdCaseTwo}a(ex)h(x)-g(x)h^{\prime}(x)=ub(x)-b(ex).
	\end{equation} 
	If we have $$\deg\big(a(ex)h(x)-g(x)h^{\prime}(x)\big)=\max\big(\deg(a(ex)h(x)),\deg(g(x)h^{\prime}(x))\big)$$ then from (\ref{equationThirdCaseTwo}) it follows in particular that $\deg(h)\leq \deg(b)$ and we are done. Otherwise, $g(x)$ is nonzero and the leading coefficients of $a(ex)h(x)$ and $g(x)h^{\prime}(x)$ cancel out. Thus, if we let $\lc(a)$ and $\lc(g)$ be respectively the leading coefficients of $a(x)$ and $g(x)$ then $\deg(h)=\lc(a)\lc(g)^{-1}$, which proves that $\deg(h)$ depends only on $\delta$ and not on the automorphism $\sigma\in\auto_{\k}(\a_{c,q},\delta)$.
	
	\medskip -- The case $f_1(x,y)=f_1(x)\in\k[x]$ and $f_2(x,y)=f_2(x)\in\k[x]$. From the assumption that $\delta$ is not locally nilpotent we deduce that $f_1(x)\neq 0$. On the other hand, using similar arguments as in the first three cases we obtain $$f_2(ex)=e^{kn}uf_2(x)+h^{\prime}(x)f_1(x)$$ and hence $\deg(h)\leq \deg(f_2)+1$.
\end{proof}

\end{document}
